\subsection{相对论与天体物理补充}

\begin{结论}{速度变换（含横向分量）}
系 $S'$ 以 $v$ 沿 $x$ 运动：
\[
\重点{u_x'=\frac{u_x-v}{1-u_xv/c^2},\qquad u_y'=\frac{u_y}{\gamma(1-u_xv/c^2)}}.
\]
快度可加 $\eta=\operatorname{artanh}\frac uc$ 只对共线速度成立；不共线 boost 的复合含转动（Thomas--Wigner）。光速不变：$u=c$ 在任何系中仍为 $c$，方向改变（光行差）。
\end{结论}

\begin{结论}{固有时：惯性世界线最长}
$\Delta\tau=\int\sqrt{1-v^2/c^2}\,\dd t$。两事件间的类时世界线中，\重点{\text{惯性（测地）世界线的固有时最长}}（双生子佯谬：转向者年轻）。
以固有加速度 $g$ 单程飞行（前一半加速、后一半减速），船时 $T$ 可达距离 $\frac{2c^2}{g}\bigl(\cosh\frac{gT}{2c}-1\bigr)$：$1g$ 下 $T=2.5$ 年约 $1.8$ 光年，$T=20$ 年约 $3\times10^4$ 光年（银河系中心量级）。
\end{结论}

\begin{结论}{逆 Compton 散射}
能量 $\epsilon$ 的光子被 Lorentz 因子 $\gamma$ 的电子散射（Thomson 极限 $\gamma\epsilon\ll m_ec^2$）：
\[
\重点{\langle\epsilon'\rangle=\frac43\gamma^2\epsilon},\qquad \epsilon'_{\max}=4\gamma^2\epsilon .
\]
单电子损能率 $P=\frac43\sigma_Tc\gamma^2\beta^2U_{\text{rad}}$（与同步辐射公式中 $U_B$ 换为 $U_{\text{rad}}$ 相同）。CMB 光子经 $\gamma\sim10^4$ 电子变成 X 射线；Sunyaev--Zel'dovich 效应。
\end{结论}

\begin{结论}{电子对产生的阈值}
光子在静止靶质量 $M$ 附近产生 $e^+e^-$：$E_\gamma\ge2m_ec^2\bigl(1+\frac{m_e}M\bigr)$——原子核附近 $\approx1.022$ MeV，\重点{\text{在电子附近 }4m_ec^2=2.044\text{ MeV}}。
两光子对撞：$E_1E_2(1-\cos\theta)\ge2m_e^2c^4$；TeV 伽马光子与红外背景光子作用而衰减（宇宙对 TeV 光子不透明的距离）。真空中单光子不能产生电子对（动量守恒）。
\end{结论}

\begin{结论}{飞行中衰变：$\pi^0\to\gamma\gamma$ 的能谱}
母粒子能量 $E$、速度 $\beta$，衰变为两个光子（静止系各向同性）：实验室中单个光子能量\重点文{均匀分布}在
\[
\Bigl[\frac{E(1-\beta)}2,\ \frac{E(1+\beta)}2\Bigr],
\]
两光子最小夹角 $\sin\frac{\theta_{\min}}2=\frac1\gamma$。一般二体衰变的子粒子能量在实验室中也是箱形分布（Doppler 线性变换 $\cos\theta^*$）。
\end{结论}

\begin{结论}{引力红移与 GPS}
弱场 $\frac{\Delta\nu}{\nu}=\frac{\Delta\Phi}{c^2}$：Pound--Rebka（$22.5$ m 塔）\重点{\frac{gh}{c^2}\approx2.5\times10^{-15}}。
GPS 卫星（$r\approx26600$ km，$v\approx3.9$ km/s）：引力使星钟每天快 $45.7\ \mu$s，运动使其慢 $7.2\ \mu$s，合计快 $38.5\ \mu$s（不修正每天定位误差 $\sim10$ km）。圆轨道上 $\frac{\dd\tau}{\dd t}=\sqrt{1-\frac{3GM}{rc^2}}$。
\end{结论}

\begin{结论}{白矮星与中子星的质量--半径}
非相对论简并电子：$R\propto M^{-1/3}$（越重越小）；电子变为相对论时压强 $\propto\rho^{4/3}$，存在极限
\[
\重点{M_{\mathrm{Ch}}\approx\frac{5.83}{\mu_e^2}M_\odot\approx1.44M_\odot}\ (\mu_e=2).
\]
白矮星 $\sim$ 地球大小；中子星 $R\approx10$--$12$ km、密度 $\sim$ 核密度，最大质量（TOV）$\approx2$--$2.3M_\odot$（观测到 $2.0M_\odot$ 以上的脉冲星）。
\end{结论}

\begin{结论}{Jeans 不稳定}
声速 $c_s$、密度 $\rho$ 的均匀气体，扰动 $\omega^2=c_s^2k^2-4\pi G\rho$：波长大于
\[
\重点{\lambda_J=\sqrt{\frac{\pi c_s^2}{G\rho}}}
\]
的扰动坍缩。Jeans 质量 $M_J\sim\rho\lambda_J^3\propto T^{3/2}\rho^{-1/2}$：冷而密的分子云碎裂成恒星；宇宙复合后 $M_J\sim10^5M_\odot$（球状星团量级）。
\end{结论}

\begin{结论}{恒星的标度律}
主序星 $L\propto M^{3.5}$（中等质量），寿命 \重点{t\propto\frac ML\propto M^{-2.5}}：太阳约 $10^{10}$ 年，$10M_\odot$ 约 $3\times10^7$ 年。$L=4\pi R^2\sigma T_{\text{eff}}^4$。
星等：$m_1-m_2=-2.5\log_{10}\frac{F_1}{F_2}$；距离模数 $m-M=5\log_{10}\frac{d}{10\text{ pc}}$。太阳 $M_V=4.83$，$T_{\text{eff}}=5772$ K。
\end{结论}

\begin{结论}{宇宙学距离}
共动距离 $d_C=c\int_0^z\frac{\dd z'}{H(z')}$（平直宇宙 $d_M=d_C$）。光度距离 $d_L=(1+z)d_M$，角直径距离 $d_A=\frac{d_M}{1+z}$：
\[
\重点{d_L=(1+z)^2d_A}\quad(\text{Etherington 互易}).
\]
表面亮度 $\propto(1+z)^{-4}$（Tolman 检验）。$d_A$ 在 $z\approx1.6$ 附近达最大（$\Lambda$CDM）——更远的星系看起来反而更大。
\end{结论}

\begin{结论}{宇宙的年龄与几个红移}
物质主导平直宇宙 $t_0=\frac2{3H_0}$；$\Lambda$CDM（$\Omega_m=0.31$）$t_0\approx0.95/H_0\approx13.8$ Gyr。$\frac1{H_0}\approx14.5$ Gyr（$H_0=67.7$）。
$T_{\text{CMB}}\propto1+z$：复合/退耦 $z\approx1090$（$T\approx3000$ K，远低于 $13.6$ eV 对应的温度，因为光子/重子比 $\sim10^9$）；物质--辐射相等 \重点{z\approx3400}；物质--暗能量相等 $z\approx0.3$。
\end{结论}

\begin{结论}{引力波：啁啾质量}
双星旋近的频率演化只依赖啁啾质量 $\mathcal M=\frac{(m_1m_2)^{3/5}}{(m_1+m_2)^{1/5}}$：
\[
\重点{\dot f=\frac{96}5\pi^{8/3}\Bigl(\frac{G\mathcal M}{c^3}\Bigr)^{5/3}f^{11/3}}.
\]
GW150914：$\mathcal M\approx30M_\odot$；合并前 $f\sim$ 数百 Hz。应变幅度 $h\sim\frac{4}{d}\frac{(G\mathcal M)^{5/3}(\pi f)^{2/3}}{c^4}$，量级 $10^{-21}$。引力波频率是轨道频率的两倍。
\end{结论}

\begin{结论}{吸积盘温度剖面}
几何薄、光学厚的稳态盘（吸积率 $\dot M$）：局部黑体温度
\[
\重点{T(r)=\Bigl[\frac{3GM\dot M}{8\pi\sigma r^3}\Bigl(1-\sqrt{\frac{r_{\text{in}}}r}\Bigr)\Bigr]^{1/4}}\propto r^{-3/4}.
\]
盘谱是多温黑体叠加，中间段 $F_\nu\propto\nu^{1/3}$。恒星级黑体盘内区 $\sim10^7$ K（X 射线），超大质量黑洞 $\sim10^5$ K（紫外）：$T_{\max}\propto M^{-1/4}$（同 Eddington 比时）。
\end{结论}

\begin{结论}{相对论聚束与视超光速}
以 $\beta$ 运动的源，辐射集中在前向 $\frac1\gamma$ 锥内；接近视线方向（夹角 $\theta$）运动的喷流，视横向速度
\[
\重点{\beta_{\text{app}}=\frac{\beta\sin\theta}{1-\beta\cos\theta}},\qquad \max\beta_{\text{app}}=\gamma\beta\ (\cos\theta=\beta).
\]
$\gamma=10$ 时可“看到” $10c$ 的横向运动。观测流量 Doppler 增强 $\delta^{3+\alpha}$，$\delta=\frac1{\gamma(1-\beta\cos\theta)}$（耀变体）。
\end{结论}

\begin{结论}{运动镜面反射：频率与能量}
光正入射到以 $\beta$ 远离的镜面：反射频率
\[
\重点{\frac{\omega'}{\omega}=\frac{1-\beta}{1+\beta}},
\]
迎面运动时取 $\beta\to-\beta$（相对论性飞镜产生高频辐射）。光子数守恒，所以反射能量按同样比例变化，镜子获得的功等于能量差——光帆的效率 $\frac{2\beta}{1+\beta}$。
\end{结论}

\begin{结论}{快速运动物体的视觉外形：Terrell 旋转}
Lorentz 收缩是“同时测量”的结果；照相（同时到达的光）还包含光行时差。远处小物体的照片看起来\重点{\text{是转过了一个角度而不是被压扁}}（Terrell--Penrose），球体始终呈圆形轮廓。
所以“看到收缩的物体”是错误的直观；收缩只在测量（坐标）意义上存在。
\end{结论}

\begin{结论}{Eddington 极限与 Bondi 吸积（量级）}
Eddington 光度 $L_E=\frac{4\pi GMm_pc}{\sigma_T}\approx1.26\times10^{31}\frac{M}{M_\odot}$ W。对应吸积率 $\dot M_E=\frac{L_E}{\eta c^2}$（$\eta\approx0.1$）。
球对称 Bondi 吸积：$\dot M\approx\pi\lambda\frac{(GM)^2\rho_\infty}{c_s^3}$（$\lambda\sim1$），\重点{\dot M\propto M^2}：黑洞越大吸积越快（“富者更富”）。吸积效率：中子星表面 $\sim0.15$，Schwarzschild 最内稳定轨道 $5.7\%$，极端 Kerr $42\%$。
\end{结论}

\begin{结论}{Kepler 运动的观测量：视向速度法与凌星法}
系外行星引起的恒星视向速度半振幅
\[
\重点{K=\Bigl(\frac{2\pi G}{P}\Bigr)^{1/3}\frac{m_p\sin i}{(M_*+m_p)^{2/3}}\frac{1}{\sqrt{1-e^2}}},
\]
木星对太阳约 $12.5$ m/s，地球约 $9$ cm/s。凌星深度 $\frac{\Delta F}F\approx\Bigl(\frac{R_p}{R_*}\Bigr)^2$：木星 $1\%$，地球 $8.4\times10^{-5}$；凌星概率 $\approx\frac{R_*}a$。
\end{结论}

\begin{结论}{潮汐锁定时标与地月系统}
潮汐使地球自转减慢、月球远离：现在月球每年远离约 $3.8$ cm，日长每世纪增加约 $1.7$ ms（含其他效应）。角动量守恒：地球自转角动量转移给月球轨道。
锁定时标 \重点{t\propto\frac{a^6}{M_{\text{伴}}^2}}（对距离极敏感）：月球早已锁定；水星 $3{:}2$ 自旋--轨道共振；冥王星--卡戎双向锁定。
\end{结论}

\参考{Taylor--Wheeler, \emph{Spacetime Physics}；Hartle, \emph{Gravity}；Carroll--Ostlie, \emph{An Introduction to Modern Astrophysics}；Rybicki--Lightman, \emph{Radiative Processes in Astrophysics}；Ryden, \emph{Introduction to Cosmology}；MIT 8.033、8.286、8.901。}
